%% kapa: exam=1 round=34 year=2023 period=1 form=A subject=economics subject_ko=경제학원론 seq=29 number=69
%% source: Q-Net 감정평가사 기출문제 (cst00309) artlSeq=5211789 file=raw/1008/5211789/2023년도 제34회 감정평가사 1차시험 1교시.pdf
%% conversion: kordoc 4.24.51 -> tools/gen1.py (auto); review=proofread-sonnet; flags=none
\begin{question}{69}
다음 모형에 관한 설명으로 옳지 않은 것은? (단, $Y^D$, $Y$, $C$, $I^D$, $r$, $G$, $T$는 각각 총수요, 국민소득, 소비, 투자수요, 이자율, 정부지출, 조세이며, $I_0$, $G_0$, $T_0$, $\alpha$, $\beta$, $\delta$는 모두 상수이다.)
\begin{framedtext}
○ $Y^D = C + I^D + G$ \par
○ $C = \alpha + \beta(Y - T)$, \quad $\alpha > 0$, \quad $0 < \beta < 1$ \par
○ $I^D = I_0 - \delta r$, \quad $I_0 > 0$, \quad $\delta > 0$ \par
○ $G = G_0$, \quad $T = T_0$ \par
○ $Y^D = Y$ \par
\end{framedtext}
\begin{choices}
  \item 모형에서 도출된 IS곡선은 우하향한다.
  \item $I_0$가 증가하면, IS곡선이 우측으로 이동한다.
  \item $G_0$가 증가하면, IS곡선이 우측으로 이동한다.
  \item $\beta$가 증가하면, IS곡선 기울기의 절댓값이 커진다.
  \item $\delta$가 증가하면, IS곡선 기울기의 절댓값이 작아진다.
\end{choices}
\end{question}
\answer{4}
